\documentclass[11pt,letterpaper]{article}
\usepackage[margin=0.8in,headheight=14pt,headsep=16pt]{geometry}
\usepackage[T1]{fontenc}
\usepackage{lmodern,amsmath,amssymb,mathtools,microtype}
\usepackage{xcolor,fancyhdr}
\definecolor{linkblue}{RGB}{25,55,90}
\usepackage[colorlinks=true,linkcolor=linkblue,citecolor=linkblue,urlcolor=linkblue]{hyperref}
\hypersetup{pdftitle={Integral images of Springer--Borel maps: introduction},
pdfsubject={Type B and D integral image formulas in the notation of Kumar--Rogers and Kumar--Xie}}
\pagestyle{fancy}\fancyhf{}
\fancyhead[L]{\small Integral images of Springer--Borel maps}
\fancyhead[R]{\small Draft introduction}
\fancyfoot[L]{\footnotesize 19 October 2026}
\fancyfoot[R]{\thepage}
\renewcommand{\headrulewidth}{0.3pt}
\setlength{\parindent}{1em}
\setlength{\parskip}{3pt}
\setlength{\emergencystretch}{2em}
\setlength{\abovedisplayskip}{7pt plus 2pt minus 2pt}
\setlength{\belowdisplayskip}{7pt plus 2pt minus 2pt}
\newcommand{\Z}{\mathbb Z}
\newcommand{\Q}{\mathbb Q}
\newcommand{\C}{\mathbb C}
\newcommand{\Rep}{\operatorname{Rep}}
\newcommand{\SO}{\operatorname{SO}}
\newcommand{\Sp}{\operatorname{Sp}}
\newcommand{\OG}{\operatorname{OG}}
\newcommand{\IG}{\operatorname{IG}}
\newcommand{\GL}{\operatorname{GL}}
\newcommand{\sat}{\operatorname{sat}}
\newcommand{\Sym}{\operatorname{Sym}}
\newcommand{\poly}{\mathrm{poly}}
\newcommand{\hbarprod}{\bar h_{1,k}}
\newcommand{\qb}[3]{\genfrac{[}{]}{0pt}{}{#1}{#2}_{#3}}
\begin{document}
\begin{center}
{\Large\bfseries Integral images of Springer--Borel maps}\\[3pt]
{\small An introduction in the notation of Parts II and III}
\end{center}
\vspace{-3pt}
\noindent\textbf{1. Introduction.}
The relation between representation rings and the cohomology of flag varieties
has a classical model in the homomorphism from polynomial representations of
\(\GL_k\) to the integral cohomology of a Grassmannian. Under this map,
irreducible representations indexed by partitions in the appropriate rectangle
map to Schubert classes. Springer morphisms extend this construction to a
representation-dependent setting for reductive groups.

Let \(G\) be a connected complex reductive group, \(B\) a Borel subgroup,
\(T\subset B\) a maximal torus, and \(P\supset B\) a parabolic subgroup with Levi
subgroup \(L\supset T\). Write \(\mathfrak g,\mathfrak t\) for the Lie algebras
and \(W_L\) for the Weyl group of \(L\).
An irreducible representation \(V(\lambda)\) with finite kernel determines
the Springer morphism \(\theta_\lambda:G\to\mathfrak g\), obtained by
trace-orthogonal projection of its matrix realization onto \(\mathfrak g\).
Its pullback on \(T\) identifies \(S(\mathfrak t^*)^{W_L}\), the invariant
polynomial functions on \(\mathfrak t\), with the subalgebra
\(\Rep^{\C}_{\lambda\text{-}\poly}(L)\) of the complexified representation ring.
The Borel characteristic map sends a weight to the first Chern class of its
associated line bundle on \(G/B\). Transporting its restriction to
\(W_L\)-invariants through this identification gives
\[
 \xi^P_\lambda:\Rep^{\C}_{\lambda\text{-}\poly}(L)
       \longrightarrow H^*(G/P,\C).
\]
This is the construction of Kumar~\cite{K}. For the defining representations
of the classical groups we abbreviate \(\lambda\text{-}\poly\) to \(\poly\).

Kumar--Rogers~\cite{KR} studied the restrictions to the symplectic and odd
orthogonal factors of the Levi subgroups. Their maps
\(\xi_{n,k}\) and \(\bar\xi_{n,k}\) have targets
\(H^*(\IG(n-k,2n),\C)\) and \(H^*(\OG(n-k,2n+1),\C)\), respectively.
They determined the generator images and proved that the induced stable
maps \(\xi_k\) and \(\bar\xi_k\) are injective
\cite[Theorems 16 and 29]{KR}. Here stability means the graded inverse
limit under the standard embeddings as \(n\to\infty\), with \(k\) fixed.
Kumar--Xie~\cite{KX} obtained the even orthogonal counterpart
\(\xi^D_{n,k}\), with target \(H^*(\OG(n-k,2n),\C)\), and proved stable
injectivity \cite[Theorem 3.8]{KX}. We retain the superscript \(D\) from
their introduction to distinguish this map from the symplectic one.

We keep the parameters \(n,k\) of these papers:
\(1\le k\le n-1\) in type B and \(n\ge4,\ 2\le k\le n-1\) in type D.
The variety \(\OG(n-k,N)\) parametrizes isotropic \((n-k)\)-planes
for a nondegenerate symmetric form on \(\C^N\).
On the torus put \(\bar t_j=(t_j-t_j^{-1})/2\).
On the rank-\(k\) orthogonal factor put \(\bar h_i=(h_i-h_i^{-1})/2\), identified
with \(\bar t_{n-k+i}\) by the Levi inclusion, and write
\[
 e_i=e_i\bigl((\bar h_1)^2,\ldots,(\bar h_k)^2\bigr),
 \qquad \hbarprod=\bar h_1\cdots\bar h_k,\qquad \hbarprod^{\,2}=e_k.
\]
Here \(e_i(\,\cdot\,)\) is the \(i\)-th elementary symmetric polynomial;
\(\hbarprod\) occurs only in type D. For either orthogonal map the published
generator formula is
\begin{equation}\label{eq:chern}
 e_i\longmapsto c_i^2+2\sum_{j=1}^{i}(-1)^j c_{i+j}c_{i-j}.
\end{equation}
We use \(c_i=c_i(Q^B)\) in type B and \(c_i=c_i(Q)\) in type D,
where the indicated bundle is the quotient of the trivial ambient bundle
by the tautological isotropic subbundle; \(c_0=1\) and \(c_i=0\) outside
its rank. These are \cite[Proposition 24]{KR} and \cite[Theorem 3.6]{KX}.

\newpage
The integral forms used in the earlier injectivity arguments are
\begin{align*}
 \Rep^\Z_{\poly}(\SO(2k+1))&=\Z[e_1,\ldots,e_k],\\
 \Rep^\Z_{\poly}(\SO(2k))&=\Z[e_1,\ldots,e_k]
                    \oplus\hbarprod\,\Z[e_1,\ldots,e_k].
\end{align*}
The second decomposition is additive; multiplication uses
\(\hbarprod^2=e_k\). These are the \emph{chosen normalized polynomial
\(\Z\)-forms} of \cite{KR,KX}. As in those papers, use cohomological
degrees \(\deg e_i=4i\) and \(\deg\hbarprod=2k\).
Equation~\eqref{eq:chern} and the integral Euler image show that the restrictions
\(\bar\xi^\Z_{n,k}\) and \(\xi^{D,\Z}_{n,k}\) take values in integral cohomology.

For \(X=\OG(n-k,2n+1)\) or \(\OG(n-k,2n)\), let \(L_{n,k}(d)\)
be the image of the degree-\(2d\) part of the corresponding normalized form.
We study the finite quotient
\begin{equation}\label{eq:saturation}
 S_{n,k}(d)/L_{n,k}(d),\qquad
 S_{n,k}(d)=
 \bigl(L_{n,k}(d)\otimes_\Z\Q\bigr)\cap H^{2d}(X,\Z).
\end{equation}
The group \(S_{n,k}(d)\) is the \emph{saturation} of the image.
Because integral cohomology has a Schubert basis, the inclusion
\(L_{n,k}(d)\subseteq S_{n,k}(d)\) has positive Smith factors
\(s_1\mid\cdots\mid s_r\): after integral basis changes, it is diagonal
with these entries. Thus
\(S_{n,k}(d)/L_{n,k}(d)\cong\bigoplus_j\Z/s_j\Z\).
We retain unit factors, recording the rank as well as the quotient.
This is an index inside the selected rational image; that image can
have smaller rank than the ambient cohomology.

Our first result determines all these factors in type B. Put
\[
 \ell=\min(k,n-k),\qquad A=\max(k,n-k),\qquad
 r=\lfloor n/2\rfloor,\qquad c=\lceil n/2\rceil.
\]
For a partition \(\gamma\subseteq A^\ell\), meaning
\(A\ge\gamma_1\ge\cdots\ge\gamma_\ell\ge0\), set
\begin{equation}\label{eq:statistic}
 b_i=\gamma_i+\ell-i,\qquad
 u(\gamma)=\#\{i:b_i<r\},\qquad
 E_{n,k}(\gamma)=\#\{i:b_i\ge c+u(\gamma)\}.
\end{equation}
Zeros are included to make \(\gamma\) have exactly \(\ell\) parts,
and \(|\gamma|=\sum_i\gamma_i\).

\medskip\noindent
\textbf{Theorem A (type B).}
\emph{For \(1\le k\le n-1\), the nonzero Smith factors of
\(\bar\xi^\Z_{n,k}\) in cohomological degree \(4m\) are the multiset}
\begin{equation}\label{eq:typeB}
 \bigl\{2^{E_{n,k}(\gamma)}:
       \gamma\subseteq A^\ell,\ |\gamma|=m\bigr\}.
\end{equation}
\emph{The image is zero in cohomological degrees not divisible by four.
In particular every factor is a power of two, with exponent at most \(\ell\).}

\medskip
The statistic labels multiplicities; it does not assert that the individual
Schur polynomials already give a Smith basis.
For the type-D statement, package these multiplicities as
\begin{equation}\label{eq:Bseries}
 \mathcal B_{n,k,e}(z)=
 \sum_{\substack{\gamma\subseteq A^\ell\\E_{n,k}(\gamma)=e}}
 z^{2|\gamma|}.
\end{equation}
Thus \(z^d\) records cohomological degree \(2d\), and \(e\ge0\) records
the factor \(2^e\). Set \(\mathcal B_{k,k,0}=1\) and
\(\mathcal B_{k,k,e}=0\) for \(e>0\); this is the point boundary needed
below. Parameters in \eqref{eq:statistic} are always recomputed for
the displayed pair \(n,k\).

The distinction between normalized and literal character lattices matters.
The chosen \(\Rep^\Z_{\poly}\) need not equal the intersection with the
ordinary integral representation ring. The alternative character lattice
generated by \(2^{2i}e_i\) and, in type D, \(2^k\hbarprod\), has
degree-\(2d\) image \(2^dL_{n,k}(d)\); each Smith factor is therefore
multiplied by \(2^d\).
The theorems concern the normalized forms specified above.

\newpage
Let \(M_D(n,k;e,d)\) be the multiplicity of \(2^e\) among the nonzero
Smith factors of \(\xi^{D,\Z}_{n,k}\) in \(H^{2d}(\OG(n-k,2n),\Z)\).
Put \(\mathcal I_{n,k}(z)=0\) for odd \(k\). For \(k=2h\), define
\begin{equation}\label{eq:overlap}
 \mathcal I_{n,2h}(z)=
 z^{\,2h(\lfloor(n-2h)/2\rfloor+1)}
 \qb{h+\lceil(n-2h)/2\rceil-1}{h}{z^2}.
\end{equation}
Here the Gaussian polynomial
\(\qb{v}{u}{q}=\sum_{\mu\subseteq(v-u)^u}q^{|\mu|}\)
counts partitions in a rectangle.

\medskip\noindent
\textbf{Theorem B (type D).}
\emph{For \(n\ge4\) and \(2\le k\le n-1\), every nonzero Smith factor of
\(\xi^{D,\Z}_{n,k}\) is a power of two, and for every \(e\ge0\),}
\begin{equation}\label{eq:typeD}
 \boxed{\displaystyle
 \sum_{d\ge0}M_D(n,k;e,d)z^d
  =\mathcal B_{n,k,e}(z)+z^k\mathcal B_{n-1,k,e}(z)
       +(\delta_{e,1}-\delta_{e,0})\mathcal I_{n,k}(z).}
\end{equation}
\emph{Here \(\delta\) is the Kronecker delta.}

\medskip
The first two terms are the factors of the square-generated image and
the image obtained by multiplication by \(\hbarprod\), respectively.
When \(k\) is odd they occupy different parities of \(d\). For even \(k\),
the sum of their separate saturations may have finite index in the
saturation of their combined image; we call this \emph{integral gluing}.
The polynomial \(\mathcal I_{n,k}\) counts every gluing direction.
Each such direction replaces a pair of unit factors \((1,1)\) by
\((1,2)\), with no further correction among the higher factors.
For example, \(n=4,k=2\) in cohomological degree \(8\) gives the full
type-D list \((1,2,2)\), and hence quotient \((\Z/2\Z)^2\).

The proofs refine the earlier generator calculations by integral
incidence and lattice arguments. In type B, Schur-product filtrations
are shown to be primitive, meaning that their quotients in the ambient
free abelian lattice are torsion free. A unit triangular minor and a
matching geometric rank bound determine these filtrations and yield
\eqref{eq:typeB}. In type D, relative orthogonal incidence bundles,
the integral Euler relation, and parity retractions establish the
comparison of the square and Euler sources with type B. An index
comparison then proves that the unit gluing exhausts all saturated
gluing, giving \eqref{eq:typeD}. The integral Schubert calculus is
published input; the finite-rank image classification is the additional
deduction developed in the companion proofs~\cite{Notes}.
Exact calculations corroborate the formulas in finite cases.
Together, the results determine the degreewise integral divisibility
of the selected Springer--Borel images, refining the explicit images
and stable injectivity of Parts II and III.

\vspace{2pt}
\noindent\textbf{References}
{\small
\renewcommand{\section}[2]{}
\begin{thebibliography}{KX24}
\setlength{\itemsep}{2pt}
\setlength{\parskip}{0pt}
\bibitem[K16]{K}
S.~Kumar, \emph{Representation ring of Levi subgroups versus cohomology
ring of flag varieties}, Math.\ Ann.\ \textbf{366} (2016), 395--415.
\href{https://kumar.math.unc.edu/papers/kumar76.pdf}{Author text}.
\bibitem[KR20]{KR}
S.~Kumar and S.~Rogers, \emph{Representation ring of Levi subgroups versus
cohomology ring of flag varieties II}, J.~Algebra \textbf{556} (2020),
340--362. \href{https://kumar.math.unc.edu/papers/kumar81.pdf}{Author text}.
\bibitem[KX24]{KX}
S.~Kumar and J.~Xie, \emph{Representation ring of Levi subgroups versus
cohomology ring of flag varieties III}, J.~Pure Appl.\ Algebra
\textbf{228} (2024), 107539.
\href{https://kumar.math.unc.edu/papers/kumar86.pdf}{Author text}.
\bibitem[N]{Notes}
\emph{Integral Springer--Borel image notes}, unpublished companion proofs,
audited 19 October 2026.
\href{https://github.com/sean-rogers-t/Sage-programs/blob/b16f582e29947a402672d160b363d8c9a2f60c11/borel/TYPE_B_GENERAL_SMITH.md}{Type B};
\href{https://github.com/sean-rogers-t/Sage-programs/blob/b16f582e29947a402672d160b363d8c9a2f60c11/borel/TYPE_D_EVEN_GLUING_THEOREM.md}{type D};
\href{https://github.com/sean-rogers-t/Sage-programs/blob/b16f582e29947a402672d160b363d8c9a2f60c11/borel/TYPE_D_EVEN_GLUING_AUDIT.md}{audit}.
The deductions are internally checked; external expert review and
originality remain unestablished.
\end{thebibliography}
}
\end{document}
